\chapter{Introduction}
\label{chap:introduction}

The Internet of Things connects constrained sensing and control devices to IP networks. IPv6 over Low-Power Wireless Personal Area Networks, known as 6LoWPAN, supports compressed IPv6 communication over constrained links, while RPL organises multi-hop communication to a border router or DODAG root \cite{rpl6550}. RPL routing decisions depend on advertised rank, control-plane messages, neighbour behaviour, and link conditions. These properties enable efficient routing but also create attack surfaces.

An attacker can disrupt forwarding, manipulate rank, flood the control plane, influence parent choice, create implausible topology, or manipulate identities visible through routing messages. Intrusion detection is therefore important. A detector that performs well when trained and tested on one known attack can still fail when the observed malicious mechanism changes. This problem is especially relevant to RPL because blackhole, sinkhole, selective-forwarding, rank, and identity-related attacks expose different evidence to an IDS.

SVELTE demonstrated that protocol-aware routing-state monitoring can support lightweight intrusion detection in RPL-based IoT networks \cite{svelte2013}. Pasikhani, Clark, and Gope argue that 6LoWPAN observations evolve over time and that fixed IDS models may degrade under concept drift \cite{pasikhani2022arl}. This dissertation does not claim to reproduce their full adversarial reinforcement-learning system. Instead, it develops a controlled Cooja-based extension that tests whether an IDS trained on one attack family generalises to another and whether limited evidence from the new environment restores held-out performance.

The project implements blackhole, sinkhole, DIS flooding, grayhole, increase-rank, DIO suppression, worst-parent, wormhole, and Sybil attacks in Contiki-NG/Cooja. Each family has five attack and five matched control runs. The resulting frozen dataset contains 810 labelled 60-second routing windows. Simulator attack markers are excluded from model inputs, and complete simulation seeds remain separate across training, adaptation, and testing.

The research questions are: can a static IDS transfer across RPL attack mechanisms; does whole-seed adaptation recover held-out performance; can routing-state evidence detect a controlled blackhole-to-sinkhole change online; how do attacker relocation and radio loss affect results; and does a direct rank-parent intervention provide a viable sinkhole defence?

This dissertation is simulation-based. Cooja provides controlled timing, matched controls, reproducible configurations, and known ground truth. The conclusions apply to the evaluated scenarios and do not claim physical deployment performance, calibrated energy use, mote ROM/RAM use, or network lifetime.

\section{Problem Statement}

Many IDS evaluations answer a narrow question: can a model distinguish known malicious samples from normal samples drawn from the same dataset? That question is useful, but it does not fully represent an evolving network. A high score can result from stable dataset-specific signatures, repeated observations from the same run, or features that encode the attack implementation. When the deployment encounters a different attack mechanism, those shortcuts may disappear.

RPL makes this problem unusually clear because its attacks are heterogeneous. A forwarding attacker changes delivery. A sinkhole may continue forwarding while advertising a false rank. A control-plane flood changes message frequency. A wormhole can forward reliably through an implausible path. A Sybil attacker changes identity observations. If a detector learns that attack means packet loss, it may be blind to a sinkhole or Sybil node. The research problem is therefore the difference between detection within a known attack distribution and generalisation after the meaning of malicious behaviour changes.

This project operationalises the problem by treating each attack family as a controlled environment. Source-family data define the original detector context. Target-family data define the changed context. The comparison preserves the simulator, extraction interval, labels, and broad network design. The approach cannot isolate every causal factor, but it makes the mechanism change explicit and repeatable.

\section{Aim and Objectives}

The aim is to evaluate and improve the robustness of an RPL IDS under controlled changes in attack mechanism and network condition. Six objectives support that aim:
\begin{enumerate}
  \item implement matched attack and control applications for a mechanism-diverse RPL campaign;
  \item generate a reproducible Cooja dataset with complete seed provenance and a known activation boundary;
  \item quantify static within-family and cross-family classifier behaviour without using attack implementation markers;
  \item measure recovery when limited complete target seeds are added for adaptation;
  \item evaluate online drift signals and robustness under attacker relocation and lossy radio;
  \item test whether rank-based sinkhole evidence can safely support a routing response.
\end{enumerate}

These objectives combine reproduction and extension. The attack campaign and supplied dataset baseline connect the work to the supervisor's paper. The cross-attack split, whole-seed adaptation protocol, Sybil surface, robustness conditions, and negative defence experiment form the principal extensions.

\section{Contributions}

The first contribution is an executable multi-attack Cooja campaign. The value is not the number of source files alone, but the common delayed-activation design and matched controls. The second contribution is the frozen 810-window dataset and audit, which make data provenance and marker exclusions explicit. The third is a 72-pair cross-attack matrix that exposes how often apparent IDS competence fails to transfer.

The fourth contribution is the adaptation curve. Rather than reporting only before and after retraining, it measures zero, one, two, and three complete target seeds across possible combinations. The fifth is the Sybil extension, which adds identity-related control-plane behaviour to the drift evaluation. The sixth is the combined online, robustness, and defence analysis. Together these experiments show that detection, adaptation, and mitigation are related but distinct problems.

\section{Scope and Boundaries}

The project focuses on network and routing evidence generated by RPL-lite in Contiki-NG. It does not evaluate application payload confidentiality, key management, firmware exploitation, or Internet-side firewall attacks. It does not claim a complete implementation of SVELTE or of the Gope adversarial RL framework. CART, Gaussian classification, CUSUM, and DDM are used to make behaviour interpretable and reproducible.

The trust-derived features and rank-parent intervention are secondary analyses. Trust scores are offline transformations of observed features. The online intervention is a narrow prototype whose failure is reported. An LLM explanation package was considered during development, but it is not included as a validated empirical result because a live, fixed-model evaluation was not completed.

\section{Dissertation Structure}

Chapter~\ref{chap:literature} reviews RPL security, lightweight IDS design, concept drift, adaptive learning, and recent resource-aware anomaly detection. Chapter~\ref{chap:methodology} defines the simulation, dataset, split, metrics, online monitors, robustness conditions, and defence evaluation. Chapter~\ref{chap:implementation} explains the attack hooks, applications, feature pipeline, models, and validation. Chapter~\ref{chap:results} reports the complete evidence. Chapter~\ref{chap:discussion} interprets the findings relative to SVELTE and the Gope framework, then defines limitations and threats to validity. Chapter~\ref{chap:conclusion} summarises the contribution and future work.
